%2multibyte Version: 5.50.0.2953 CodePage: 1252

\documentclass{article}
%%%%%%%%%%%%%%%%%%%%%%%%%%%%%%%%%%%%%%%%%%%%%%%%%%%%%%%%%%%%%%%%%%%%%%%%%%%%%%%%%%%%%%%%%%%%%%%%%%%%%%%%%%%%%%%%%%%%%%%%%%%%%%%%%%%%%%%%%%%%%%%%%%%%%%%%%%%%%%%%%%%%%%%%%%%%%%%%%%%%%%%%%%%%%%%%%%%%%%%%%%%%%%%%%%%%%%%%%%%%%%%%%%%%%%%%%%%%%%%%%%%%%%%%%%%%
\usepackage{amsfonts}
\usepackage{amsmath}

\setcounter{MaxMatrixCols}{10}
%TCIDATA{OutputFilter=LATEX.DLL}
%TCIDATA{Version=5.50.0.2953}
%TCIDATA{Codepage=1252}
%TCIDATA{<META NAME="SaveForMode" CONTENT="1">}
%TCIDATA{BibliographyScheme=Manual}
%TCIDATA{Created=Tuesday, September 13, 2016 11:26:35}
%TCIDATA{LastRevised=Monday, December 13, 2021 21:39:58}
%TCIDATA{<META NAME="GraphicsSave" CONTENT="32">}
%TCIDATA{<META NAME="DocumentShell" CONTENT="Standard LaTeX\Blank - Standard LaTeX Article">}
%TCIDATA{Language=American English}
%TCIDATA{CSTFile=40 LaTeX article.cst}

\newtheorem{theorem}{Theorem}
\newtheorem{acknowledgement}[theorem]{Acknowledgement}
\newtheorem{algorithm}[theorem]{Algorithm}
\newtheorem{axiom}[theorem]{Axiom}
\newtheorem{case}[theorem]{Case}
\newtheorem{claim}[theorem]{Claim}
\newtheorem{conclusion}[theorem]{Conclusion}
\newtheorem{condition}[theorem]{Condition}
\newtheorem{conjecture}[theorem]{Conjecture}
\newtheorem{corollary}[theorem]{Corollary}
\newtheorem{criterion}[theorem]{Criterion}
\newtheorem{definition}[theorem]{Definition}
\newtheorem{example}[theorem]{Example}
\newtheorem{exercise}[theorem]{Exercise}
\newtheorem{lemma}[theorem]{Lemma}
\newtheorem{notation}[theorem]{Notation}
\newtheorem{problem}[theorem]{Problem}
\newtheorem{proposition}[theorem]{Proposition}
\newtheorem{remark}[theorem]{Remark}
\newtheorem{solution}[theorem]{Solution}
\newtheorem{summary}[theorem]{Summary}
\newenvironment{proof}[1][Proof]{\noindent\textbf{#1.} }{\ \rule{0.5em}{0.5em}}
\setlength{\textheight}{9.3in}
\setlength{\topmargin}{-0.8in}
\if@twoside
  \setlength{\oddsidemargin}{0.10in}
  \setlength{\evensidemargin}{-0.30in}
\else
  \setlength{\oddsidemargin}{-0.30in}
  \setlength{\evensidemargin}{0.00in}
\fi
\setlength{\textwidth}{7.1in}
\newcommand{\m}[3]{\multicolumn{#1}{#2}{#3}}
\newcommand{\ep}{\varepsilon}
\def \dsp{\def \baselinestretch{1.1}\large \normalsize}
\dsp
\input{tcilatex}
\renewenvironment{abstract}
 {\small
  \begin{center}
  \bfseries \abstractname\vspace{-0.5em}\vspace{0pt}
  \end{center}
  \list{}{    \setlength{\leftmargin}{0mm}    \setlength{\rightmargin}{\leftmargin}  }  \item\relax}
 {\endlist}
\begin{document}


\begin{center}
{\Large Final Exam}

\bigskip 
\end{center}

\textbf{Question 1}

(a) Provide a definition for the bias of an estimator. How is the expression
for the Mean Squared Error (MSE)\ of an estimator related to its bias?

(b) Provide a definition for an efficient estimator. Explain why this is a
desirable property and how it is related to the Mean Squared Error (MSE)\ of
an estimator.

(c) Define the score vector and the Fisher information matrix.

(d) State and prove the Cramer Rao theorem for a one-dimensional parameter $%
\theta $.

\bigskip 

\textbf{Question 2}

Consider the linear regression model $y=X\beta +\varepsilon ,$ where $y$ is
an $n\times 1$ vector of dependent variables, $X$ is an $n\times k$ matrix
of regressors, $\beta $ is a $k\times 1$ vector of unknown parameters, and $%
\varepsilon $ is an $n\times 1$ vector of unobserved random variables. The
standard assumptions provided in the lecture are assumed to hold.

(a) Suppose there are two explanatory variables, $k=2,$ so that $%
y=X_{1}\beta _{1}+X_{2}\beta _{2}+\varepsilon .$ State without proof the
Frisch-Waugh-Lovell theorem for the OLS estimator $\hat{\beta}_{1}.$

(b) Consider two economists estimating the following regression models:
Economist 1 estimates $y=X_{1}\beta _{1}+\varepsilon $ while Economist 2
also includes $X_{2}$ and estimates instead $y=X_{1}\beta _{1}+X_{2}\beta
_{2}+\varepsilon .$ What are the estimators of the two economists for $\beta
_{1}$?

(c) Given that the true data generating process is $y=X_{1}\beta
_{1}+X_{2}\beta _{2}+\varepsilon ,$ show that that the estimator for $\beta
_{1}$ of Economist 2 is unbiased.

(d) Given that the true data generating process is $y=X_{1}\beta
_{1}+X_{2}\beta _{2}+\varepsilon ,$ show that the estimator for $\beta _{1}$
of Economist 1 is biased. Discuss this result in relation to omitted
variable bias.

(e)\ Now suppose the true data generating process is $y=X_{1}\beta
_{1}+\varepsilon ,$ so $\beta _{2}=0.$ Are the two estimators of the two
economists unbiased?

(f) Given that the true $\beta _{2}=0,$ compare the variance of $\hat{\beta}%
_{1}$ for the two economists. What can you conclude about the tradeoff
between bias and variance when deciding which explanatory variables to
include?

\bigskip 

\textbf{Question 3 }

Consider the Gaussian linear model $y=X\beta +\varepsilon ,$ where $\beta
\in \mathbb{R}^{K}$ and $\varepsilon |X\sim \mathcal{N}\left( 0,\sigma
^{2}\Omega \right) $ for an unknown scalar $\sigma ^{2}$ and some \textbf{%
known} $n\times n$ positive definite matrix $\Omega $.

(a) Write down the log-likelihood function of the sample. 

HINT: Recall the definition of a Normal density of an $m\times 1$ vector $%
z\sim \mathcal{N}\left( \mu ,\Sigma \right) $: $\ $%
\begin{equation*}
f\left( z;\mu ,\Sigma \right) =\left( 2\pi \right) ^{-m/2}\det \left( \Sigma
\right) ^{-1/2}\exp \left\{ -\frac{1}{2}\left( z-\mu \right) ^{\prime
}\Sigma ^{-1}\left( z-\mu \right) \right\} .
\end{equation*}

(b) Show that the maximum likelihood estimator $\hat{\beta}^{_{MLE}}$ of $%
\beta $ coincides with the generalised least squares estimator $\hat{\beta}%
^{_{GLS}}$ of $\beta $.

(c) Show that $\hat{\beta}^{_{MLE}}$ is unbiased and calculate its variance
conditional on $X$.

(d) Find the maximum likelihood estimator for $\sigma ^{2}$.

(e) Compute the Fisher Information matrix of $[\beta ,\sigma ^{2}]$
conditional on $X$.

(f) Does the $\hat{\beta}^{_{MLE}}$ achieves the Cramer Rao Lower bound
[conditional on $X$]? How does this relate to the Gauss Markov theorem for
the GLS estimator?

\end{document}
